\section{Introduction}
\label{sec:introduction}

Modern industrial and telecommunication infrastructures increasingly rely on distributed Service-Oriented Architectures (SOA) operating across geo-distributed cloud environments and edge compute nodes. While microservice decoupling offers modularity, independent scalability, and deployment flexibility, it fundamentally expands the system's attack surface~\citep{saltzer1975protection}. In large-scale telemetry networks, such as smart city sensor grids, environmental monitoring systems, and critical infrastructure telemetry, thousands of field sensors transmit continuous readings across untrusted Wide Area Networks (WAN) to centralized ingestion gateways, storage nodes, and real-time analytic processing engines.

Deploying such systems at scale presents significant architectural and security challenges:
\begin{enumerate}
    \item \textbf{Dynamic Scalability and Namespace Chaos:} As microservice instances are dynamically scaled across physical virtualization hosts and Linux container runtimes (such as Incus/LXD), flat service discovery registries degrade. Locating which physical hardware or virtual container hosts an anomalous service instance becomes computationally prohibitive without a structured, hierarchical naming convention.
    \item \textbf{High-Volume Adversarial Vulnerability:} Unprotected ingestion gateways are exposed to hostile network environments, making them prime targets for command injection attacks (e.g., Shellshock CVE-2014-6271~\citep{cve20146271}), volumetric Denial-of-Service (DDoS) flood botnets, and in-flight Man-in-the-Middle (MitM) telemetry tampering.
    \item \textbf{The Observability-Security Gap:} Operational monitoring tools (e.g., Prometheus~\citep{prometheus2026}) track metric regressions (latency surges, elevated error rates, CPU starvation) but remain blind to malicious intent. Conversely, Security Information and Event Management (SIEM) systems (e.g., Wazuh~\citep{wazuh2026}) detect attack signatures in access logs but lack contextual awareness of operational resource degradation. Without automated correlation between these operational and security silos, incident localization and root cause analysis remain manual, slow, and error-prone.
    \item \textbf{WAN Bandwidth Inefficiencies:} In multi-region deployments, naïve global polling across WAN links introduces severe bandwidth consumption and monetary transit costs.
\end{enumerate}

To systematically resolve these challenges, this project presents the end-to-end design, implementation, formal threat modeling, and quantitative evaluation of a hardened, multi-tier SOA telemetry architecture designed for large-scale operation. 

\subsection{System Overview and Implemented Scope}
To satisfy the requirements of this project, I built the entire telemetry and security architecture from the ground up, including all core microservices, custom Prometheus exporters, adversary simulation tools, and automated container deployment scripts. The implemented system centers on four primary areas:
\begin{itemize}
    \item \textbf{Hierarchical Service Discovery and Incident Localization:} A dedicated Go Service Registry implementing a Fully Qualified Service Name (FQSN) schema that encodes geographic regions, network zones, physical hosts, and container boundaries, paired with an $O(1)$ incident localization API that directly resolves affected containers and hosts during alerts.
    \item \textbf{Three Syllabus Threat Defenses and Active Responses:} End-to-end defenses against three syllabus-mandated threat vectors: (1) Shellshock RCE detection and threat-intelligence IP blocking via Wazuh SIEM, (2) volumetric DoS throttling using an in-memory token-bucket algorithm and payload clamping, and (3) Route and Signals hijacking mitigation using HMAC-SHA256 integrity signatures and timestamp freshness validation.
    \item \textbf{Autonomous FCAPS Fault Management and Event Correlation:} An autonomous 5-stage closed-loop fault management engine~\citep{itu_fcaps} that ingests Prometheus operational metrics, correlates them with Wazuh access logs to deduce root causes, executes automated active responses (\texttt{firewall-drop}), and verifies baseline recovery.
    \item \textbf{Locality-Aware Scalability and WAN Bandwidth Optimization:} A locality-aware discovery mechanism evaluated across up to $N = 1000$ services, demonstrating sub-millisecond lookup latency, a $99.99\%$ cross-region WAN bandwidth reduction, and empirical profiling of storage lock contention under heavy concurrency.
\end{itemize}

\subsection{Document Outline}
The remainder of this report is organized as follows: \cref{sec:state_of_the_art} surveys the relevant state of the art in SOA, SIEM, Prometheus, and Attack-Defense modeling. \cref{sec:architecture} details the implemented multi-tier architecture and hierarchical discovery model. \cref{sec:risk_management} analyzes attack surfaces and develops formal Attack-Defense Trees. \cref{sec:three_threats} presents the detailed threat evaluation, exploitation mechanics, and mitigations for the three syllabus threats. \cref{sec:fcaps} presents the autonomous FCAPS correlation engine. \cref{sec:evaluation} reports the quantitative scalability evaluation, bandwidth proofs, and bottleneck analysis. Finally, \cref{sec:conclusion} concludes the paper with lessons learned and outlines future work.
